\section{Related work}%
\label{sec:related}

\paragraph{Optical flow.}

Optical flow approximates dense instantaneous motion.
Originally approached by studying the brightness constancy equation~\cite{horn1981determining, lucas1981iterative, Black:ICCV:1993, bruhn2005lucas}, starting with
FlowNet~\cite{dosovitskiy2015flownet,ilg2017flownet} optical flow has been tackled using deep learning.
More recently, DCFlow~\cite{Xu_2017_CVPR} has introduced the computation of a 4D cost volume, used in most follow-up works~\cite{sun2018pwc, wang2020displacement, raft}.
Notable is RAFT~\cite{raft}, which introduced incremental flow updates and inspired several follow-up works~\cite{jiang2021learning, xu2021high, jiang2021learning, zhang2021separable}.
\method also uses 4D cost volumes and iterated updates but is applied to tracking.

Transformers~\cite{vaswani2017attention} have also been applied to the optical flow problem~\cite{huang2022flowformer, shi2023flowformer++, zhao2022global}.
Flowformer~\cite{huang2022flowformer} drew inspiration from RAFT and proposed a transformer-based approach that tokenizes the 4D cost volume.
GMFlow~\cite{zhao2022global} replaced the update network with a softmax with self-attention for refinement.
Perceiver IO~\cite{jaegle22perceiver} proposed a unified transformer for several tasks, including optical flow.

Optical flow can be applied to tracking by integrating predictions over time, but this approach accumulates errors in the form of drift~\cite{pips}, which motivates developing architectures like \method that track points over longer horizons.

\paragraph{Multi-frame optical flow.}

Several authors~\cite{ren2019fusion, shi2023videoflow, janai2018unsupervised, luo2023continuous} have extended optical flow to multiple frames.
Early methods used Kalman filtering~\cite{chin1994probabilistic, elad1998recursive} to encourage temporal consistency.
Modern multi-frame methods produce dense flow.
RAFT~\cite{raft} can be applied in a warm-start manner~\cite{sui2022craft, sun2022skflow} for multi-frame optical flow estimation.
However, these methods are not designed for long-term tracking and do not consider points occluded for a long time.
VideoFlow~\cite{shi2023videoflow} extends optical flow to three and five consecutive frames by integrating forward and backward motion features during flow refinement.
MFT~\cite{neoral2023mft} conducts optical flow estimation between distant frames and chooses the most reliable chain of optical flows.
Our tracker scales to produce semi-dense tracks.

\paragraph{Visual object tracking.}

Before the emergence of deep learning, some authors proposed handcrafted joint trackers~\cite{birchfield08joint,sand2008particle}, but few have considered doing so using deep networks like \method.
Our work is weakly related to \emph{multiple object tracking}~\cite{chen22visual} where tracking through occlusions~\cite{zhao2004tracking}, with changes in appearance~\cite{matthews2004template}, and with temporal priors~\cite{sidenbladh2000stochastic} have been extensively studied.

Some recent point trackers have been inspired by visual object tracking that predict the target states by extracting and combining features extracted from source and target image regions~\cite{bertinetto2016fully, danelljan2017eco, li2018learning, Bhat_2019_ICCV, danelljan2019atom}.
Transformers have also been applied in visual object tracking~\cite{Chen_2021_CVPR, cui2022mixformer}.
However, our focus is the tracking of points, including background points, not objects.

\paragraph{Tracking any point.}

Particle Video~\cite{sand2008particle} introduced the problem of Tracking Any Point in a video (TAP), and PIPs~\cite{pips} proposed a better model which can track through occlusions more reliably by predicting tracks in a sliding window and restarting them from the last frame where the point was visible.
However, PIP cannot track a point beyond the duration of a window.
TAP-Vid~\cite{tapvid} introduced a new benchmark and a simple baseline for TAP\@.
While the original Particle Video does track points jointly, PIPs and TAP-Vid track points independently (in parallel).
Inspired by TAP-Vid~\cite{tapvid} and PIPs~\cite{pips}, TAPIR~\cite{doersch2023tapir} further improved tracking with a two-stage, feed-forward point tracker that uses TAP-Vid-like matching and PIPs-like refinement.

PointOdyssey~\cite{zheng2023pointodyssey} addressed long-term tracking with PIPs++, a simplified version of PIPs, and introduced another benchmark for long-term tracking.
However, PIPs++ still tracks points independently.
OmniMotion~\cite{wang2023tracking} optimizes a volumetric representation for each video, refining correspondences in a canonical space.
However, this approach requires test-time optimization, which, due to its cost, is not suitable for many practical applications, especially online ones.

Trackers and optical flow models are often trained using synthetic datasets~\cite{karaev2023dynamicstereo, mayer2016large, dosovitskiy2015flownet, sun2021autoflow, butler2012naturalistic, zheng2023pointodyssey}, as annotating real data can be challenging.
Synthetic datasets provide accurate annotations, and training on them has demonstrated the ability to generalize to real-world data~\cite{mayer2016large, sun2021autoflow}.


